Enfin, pour la vitesse tout au long de la glissade, les équation de la
conservation de l'énergie sont utilisé ce qui nous permet de décrire notre
système comme suit:

\begin{align}
\underbrace{mgy_0}_{\textrm{énergie totale}} =
  \underbrace{mgF(x)}_{\textrm{énergie potentiel à } x} +
	\underbrace{\frac{1}{2}m[v(x)]^2}_{\textrm{énergie cinétique à }x} +
	\underbrace{\int_0^x{F_f(x)x} \d{x}}_{\textrm{travail cummulatif de la friction
	à }x}
\end{align}

Il est intéressant de remarquer que la pente de la trajectoire (dérivée) trouvée
précédement permet d'établir un angle en radiant de la pente selon $x$ comme
suit:
\begin{equation}
	\theta(x) = \arctan{\left( F^\prime(x) \right)}
\end{equation}

Cela permet d'exprimer la force de frottement selon x (celle-ci étant définie
comme $\mu_f mg\cos{\theta}$). Puis, en isolant la vitesse nous arrivons à
l'équation suivante:

\begin{equation}
	v(x) = \sqrt{
	  2g\left( y_0 - F(x) - \mu_f \int_0^x{cos{(\arctan{(F^\prime(x))})}x \dx} \right)
	}
\end{equation}

En utilisant MATLAB pour re-calculer l'intégrale à chaque valeur de $x$ le long
de la trajectoire, il est possible de tracer plusieur courbes en posant des
des $\mu_f$ réparties dans les plage que la valve permet.













